\documentclass[11pt]{book}
\input{../../../Manuals/Bibliography/commoncommands}

\newcommand{\fdsversion}{6.11.0}
\newcommand{\textct}[1]{\texttt{#1}}
\bibliography{../../../Manuals/Bibliography/FDS_refs,../../../Manuals/Bibliography/FDS_general,../../../Manuals/Bibliography/FDS_mathcomp,scripts/FTMI_refs}

\begin{document}


\pagestyle{empty}

\begin{minipage}[t][9in][s]{6.5in}

\headerB{FDS2FTMI User's Guide}
\begin{flushright}
\fontsize{20}{24}\selectfont
\bf{
An automated code to one-way coupling \\
between FDS and FEM using FTMI \\
}
\end{flushright}
\begin{flushright}
\fontsize{14}{16}\selectfont
{
Julio Cesar Silva  
}

\end{flushright}

\headerC{
\today \\
FDS Version \fdsversion \\
Revision: 

}

\end{minipage}

\newpage

\newpage

\frontmatter

\pagestyle{plain}

\cleardoublepage
\phantomsection
\tableofcontents

\mainmatter


\chapter{Introduction}
\label{info:intro}

This document explains the developed strategy to evaluate the behavior of structures under fire conditions by a coupled analysis including Fire Dynamics Simulator (FDS\cite{FDS_Users_Guide}) and Finite Element Method (FEM) codes. Common practice is to extract the results from a fire simulation, and apply those as boundary conditions at the exposed surfaces for the thermomechanical analysis. This kind of approach is commonly referred as one-way coupling. In a two-way coupling strategy, the thermomechanical results, e.g., displacements, collapses, etc. are transposed back to the fire simulation. The two-way approach can lead to a more complex simulation, increasing the amount of data to be transferred between the models. Its advantages are related to cases where displacements can change the fire source or the fluid flow pattern, creating a different fire scenario. Using one-way coupling it is possible to establish each model separately, which means that the same fire simulation results can be used on slightly different thermomechanical models, contributing to development of optimized structures.  

The Fire-ThermoMechanical Interface \cite{FTMI:Methodology} is a one-way coupling procedure that uses the Adiabatic Surface Temperature ($T_{\hbox{\tiny AST}}$) concept \cite{Wickstrom:Interflam2007,Wickstrom_AST} to establish an interface between FDS and FEM codes. The procedure recreates the {\em net} heat flux evaluated in FDS, in the FEM environment by transferring $T_{\hbox{\tiny AST}}$ and the heat transfer coefficient ($h$) from FDS results and adopting the wall temperature ($T_{\rm w}$) provided by the FEM code at each time step  ($h$ can also be adopted as a constant). This heat flux can be calculated by FEM over the elements faces and applied into each node of the exposed surface. FDS2FTMI is a code created to perform these tasks in an automated way, allowing the simulation of the behavior of global structures, discretized with shell and/or solid elements, exposed to fire conditions. FTMI can be used with any FEM code, at this time, FDS2FTMI is implemented between FDS and A{\footnotesize NSYS} ~\footnote{Certain commercial entities, equipment, or materials may be identified in this document in order to describe an experimental procedure or concept adequately. Such identification is not intended to imply recommendation or endorsement by the author or the National Institute of Standards and Technology, nor is it intended to imply that the entities, materials, or equipment are necessarily the best available for the purpose.}.

The code FDS2FTMI is included in the FDS-SVM repository \cite{FDS-SMV_repository} under \path{Utilities/Structural_Interaction/fds2ftmi}. A makefile is provided to help in the compilation process with different operational systems and compilers. A shell and a bat scripts (\path{scripts/build_fds2ftmi.sh}, \path{scripts/build_fds2ftmi.bat}) are also provided to build this manual, run the Verification cases and check the results~\footnote{Two environment variables are needed to run these scripts: FDS\_GITROOT, which is the path for fds-smv git folder; and ANSYS, which is the path to the ANSYS executable (including the executable).}.  

%If a reference to this methodology is needed, Silva~{\textit {et. al}} \cite{FTMI:Methodology} should be used. 
Some validation work done with this methodology is provided at Zhang~{\textit {et. al}} \cite{FTMI:Validation}.

\chapter{Getting Started}
\label{info:start}

Heat can be transferred from flames and hot gases to structures surfaces by radiation and convection. The total heat flux ($\dot{q}_{\rm tot}''$) is defined by the sum of these two parcels: 
\be \label{eq_qtot} \dot{q}_{\rm tot}'' = \dot{q}_{\rm r}'' + \dot{q}_{\rm c}'' = \epsilon \, \left( {e}_{\rm inc} - \sigma T_{\rm w}^4 \right) + h (T_{\rm g} - T_{\rm w})  \ee  
where $\dot{q}_{\rm r}''$ is the radiative heat flux, $\dot{q}_{\rm c}''$ is the convective heat flux, $\epsilon$ is the emissivity, ${e}_{\rm inc}$ is the radiative thermal energy incident on the exposed surfaces, $\sigma$ is the Stefan-Boltzmann constant, $T_{\rm w}$ is the wall temperature and $T_{\rm g}$ is the gas temperature.

Advanced fire simulation models provide results that characterize the three dimensional evolution of the fire, the radiative thermal energy incident on the exposed surfaces and the gas temperatures. However, the models are not capable of accurately characterizing the temperature distribution on solids, and thus the surface temperature. Consequently, the total heat flux, as presented on Eq.~\ref{eq_qtot}, cannot be precisely calculated at the end of the fire simulation. Therefore, an additional approach is necessary to cross this barrier.

\section{Adiabatic Surface Temperature}

In order to establish an accurate interface between the fire simulation and the thermomechanical analysis, the Adiabatic Surface Temperature concept is used \cite{Wickstrom:Interflam2007}. A real surface can be replaced by a perfectly insulated surface exposed to the same heating conditions; the total heat flux to this ideal surface is by definition zero :
\be \label{eq_ast} \epsilon \, \left( {e}_{\rm inc} - \sigma T_{\hbox{\tiny AST}}^4 \right) + h (T_{\rm g} - T_{\hbox{\tiny AST}}) = 0  \ee 
and the temperature of this idealized surface can be obtained as an output from the fire simulation.

Since the real and the hypothetical surfaces are exposed to the same heating conditions, the total heat flux that will be applied to the thermomechanical model can be obtained by the solution of the system formed by Eqs. (\ref{eq_qtot}) and (\ref{eq_ast}):
\be \label{eq_ftmi}\dot{q}_{\rm tot}'' = \epsilon \sigma \, \left( T_{\hbox{\tiny AST}}^4 - T_{\rm w}^4 \right) + h (T_{\hbox{\tiny AST}} - T_{\rm w})  \ee

This single variable ($T_{\hbox{\tiny AST}}$) is considered capable of reducing the complexity of the fire simulation into one simple scalar~\cite{FTMI:Duthinh,FTMI:Sandstrom,FTMI:Wickstrom2010}. To achieve a correct definition of the total heat flux, accounting correctly both the convective and radiative heat fluxes, it is also necessary to add the convective heat transfer coefficient (\textit{h}). 
 
During the FDS simulation, $T_{\hbox{\tiny AST}}$ and $h$ can be recorded via \textct{DEVC}, \textct{BNDF} (for traditional obstructions), or \textct{GEOM} boundary outputs. To use \textct{DEVC}s, the points (or nodes) where the information will be collected need to be determined in advance, before the fire simulation. On the other hand, \textct{BNDF} and \textct{GEOM} boundary outputs generate data at all exposed surfaces (boundary cells or triangulated faces marking the interface between gas and solid) within the domain. FDS2FTMI extracts information stored in \textct{BNDF} files as well as unstructured geometry boundary files (\textct{.be} and \textct{.gcf}), meaning the coupling procedure can be achieved for different discretization levels (at FEM) and arbitrarily complex shapes. Small modifications or dimensioning in FEM do not imply restarting the fire simulation. 
  
\section{Capturing the Mesh Information from FEM}

In FDS, the domain is divided into orthogonal cells in a rectilinear grid. Traditionally, geometries are approximated by blocks (\textct{OBST}s) that snap to this grid. However, FDS also supports complex unstructured geometries (\textct{GEOM}s) represented by triangulated surfaces, which are not constrained to the rectilinear grid. For FEM thermomechanical analysis, the structure is usually discretized with solid (tetrahedrons or hexahedrons) or shell (3D plane) elements. Also, the element size needed at the thermomechanical model can lead to an unfeasible fire simulation. 

Even for modest structures, the models generated for each side of this interface may not match perfectly. Therefore, the communication between the two models is based on the spatial position of the nodes on the exposed surfaces (FEM) and the cell faces or triangulated geometry centroids (FDS). The geometries created for both models need to share a common Cartesian coordinate system, i.e., they need to have the same origin and be consistent. 

FDS2FTMI's main function is to read the \textct{BNDF} files (for \textct{OBST}s) as well as \textct{.be} and \textct{.gcf} files (for \textct{GEOM}s) to extract $T_{\hbox{\tiny AST}}$ and $h$ at the exposed surfaces. To handle complex shapes, the code evaluates closest-face centroid proximity to accurately map FDS boundary data onto the FEM nodes. This code also generates the files that need to be read by the FEM code to apply correct boundary conditions. The source code can be found at \path{Utilities/Structural_Interaction/fds2ftmi/scripts}.

The input parameters for FDS2FTMI are the spatial position of the exposed surfaces at FEM model. These parameters are divided in 2 files, ``nodes.dat'' and ``elements.dat'', and include the node positions (Cartesian coordinates) and the node numbers at the exposed faces ordered using right-hand rule with normal pointing outwards. 

\paragraph{FTMI Format (New)}~\\
With the updated Python tool, a streamlined CSV-based format is available when using the \texttt{code:ftmi} token. This format requires only a single ``nodes.dat'' file containing the target node positions and orientations.

\subparagraph{nodes.dat (FTMI Format)}~\\
The first line contains the total number of nodes. All subsequent lines are comma-separated values (CSV) detailing the node ID, X, Y, and Z Cartesian coordinates, and an optional orientation flag (1 for +x, -1 for -x, 2 for +y, -2 for -y, 3 for +z, -3 for -z). If the orientation is left blank, the script will map the closest geometry face regardless of orientation.
\begin{verbatim}
   3
   Node_1, 1.000, 0.000, 0.000, 1
   Node_2, 1.000, 0.100, 0.000, 1
   Node_3, 1.000, 0.200, 0.000, 1
\end{verbatim}

\paragraph{A{\footnotesize NSYS}}~\\ 
If using \texttt{code:ansys}, the traditional format applies. The interface with A{\footnotesize NSYS} was developed using surface effect elements, so the information is presented in two files: ``nodes.dat'' and ``elements.dat''. To help in the generation of the surface effect elements and these two input files, A{\footnotesize NSYS} macros (e.g., ``ansys2ftmi.geom'' or ``ansys2ftmi\_shell.geom'') are provided in the scripts folder.

\subparagraph{nodes.dat and elements.dat}~\\
The first line of ``nodes.dat'' has the number of nodes and the highest node number in the FEM model. All other lines have just the node number and the Cartesian coordinates (x, y and z).
\begin{verbatim}
   100   2400
   1    1.000    0.000    0.000
   2    1.000    0.100    0.000
   .
   .
   .
   100  1.000    1.000    1.000
\end{verbatim}

The first line of ``elements.dat'' has the number of interface elements that compose the exposed surface, and 3 additional flags. The first flag dictates the type of element used in the FEM model: 0 means solid elements and 1 means shell elements. The second flag defines in which layer the heat flux is prescribed: 1 means top layer and 2 means bottom layer. The last flag defines the number of the interface element type in A{\footnotesize NSYS}.
\begin{verbatim}
   200     0     0     2
   801   120     2   148   117   121   149   139   118
   802   126   128   130   151   127   144   125   150
   .
   .
   .
  1000   178   163    41    69   179   164    10    80
  END
\end{verbatim}
In this example, there are 200 interface elements in this model, they are attached to solid elements (so the first and second flag are 0), and they correspond to the element type 2 in the FEM model.

\section{Executing the code}

After the definition of the files \path{nodes.dat} and \path{elements.dat}, the Python script \path{pyfds2ftmi.py} is executed to extract information from FDS boundary files. This code accepts parameters as flexible flag-style tokens. The required and optional tokens are:

\paragraph{chid:} 
Drives the naming convention of FDS output files. It is needed for the code to get information from the \path{.smv}, \path{.bf}, \path{.be}, and \path{.gcf} files.

\paragraph{cell:} 
This is the cell size used to create a virtual search radius over the FEM node positions to map results from FDS BNDF and GEOM files. It is generally adopted as the same value as the cell size in the FDS simulation.

\paragraph{var:} 
Specifies which variables to extract. If defined as \texttt{2} or \texttt{ast+h}, the code will extract both $T_{\hbox{\tiny AST}}$ and $h$ from the FDS simulation. If defined as \texttt{1} or \texttt{ast}, the code will extract only $T_{\hbox{\tiny AST}}$. If a constant convective heat transfer coefficient is desired, it can be prescribed using the format \texttt{1:h:<value>} (e.g., \texttt{1:h:35}).

\paragraph{time:} 
(Optional) Specifies a time window for the extraction using the format \texttt{t0:t1} (e.g., \texttt{time:0:600}).

\paragraph{avg:} 
(Optional) Defines the averaging behavior. Options include \texttt{none} for no averaging, or \texttt{mean:N} to average over the last $N$ time steps.

\paragraph{code:}
Tells FDS2FTMI which FEM code format to output. Supported options include \texttt{ansys}, \texttt{ftmi}, \texttt{abaqus}, \texttt{opensees}, and \texttt{custom}.

\paragraph{file:}
The path and filename for the generated output file, which will contain the commands or data tables to be read by the thermomechanical analysis code.
\\

These tokens can be provided in any order via the command line. For example:
\begin{verbatim}
  python pyfds2ftmi.py chid:simple_panel_hot cell:0.05 var:2 time:0:600 
  avg:none code:ansys file:simple_panel.dat
\end{verbatim}

The code will search into the boundary files for the results. Searching over every surface in the FDS simulation can slow down this process. To reduce the size of the \textct{BNDF} files, add \textct{BNDF\_DEFAULT=.FALSE.} in the \textct{MISC} line and then add \textct{BNDF\_OBST=.TRUE.} at the exposed surfaces' \textct{OBST} lines.

\chapter{Verification Cases}
\label{info:verification}

This chapter contains the Verification cases that test the FDS2FTMI code. The version hash on the figures is related to the FDS repository \cite{FDS-SMV_repository} Git commit hash used to compile the codes and run the cases.  

\section{Simple panel exposed to radiation and convection (\texorpdfstring{\textct{simple\_panel\_hot}}{simple\_panel\_hot})}

This example verifies the FDS2FTMI ability to translate the heat transfer between solid and gas phase at fire simulation into boundary conditions for a FEM model (A{\footnotesize NSYS}). It also demonstrates that the code can extract correctly the results from \textct{BNDF} files. This model is composed by 2 panels, one has a constant temperature of $T_s=1000$~\si{\degree C} and the other one has an initial temperature of $T_s=20$~\si{\degree C}. The temperature evolution at the later panel is the target of this exercise (the input files are provided in \path{examples/simple_panel_hot}). The case is set to be 2D, with only one cell at $y$ axis. These panels are $1$~m height ($z$ axis) and they are $1$~m from each other. Three \textct{DEVC}s were setup to compare the results, they are placed at $0.775$~m (1), $0.525$~m (2) and $0.275$~m (3) in $z$ axis. After the fire simulation, the boundary conditions are translated to A{\footnotesize NSYS} by FTMI. A comparison between the Adiabatic Surface Temperature $T_{\hbox{\tiny AST}}$ and the surface temperature $T_s$ obtained by FTMI and FDS is provided in Fig.~\ref{simple_panel_hot}. 

\begin{figure}[ht]
\noindent
\begin{tabular*}{\textwidth}{l@{\extracolsep{\fill}}r}
\includegraphics[width=3.0in]{scripts/SCRIPT_FIGURES/simple_panel_hot_ast} &
\includegraphics[width=3.0in]{scripts/SCRIPT_FIGURES/simple_panel_hot_ts}
\end{tabular*}
\caption[The \textct{simple\_panel\_hot} results]{Comparison of the results obtained by FTMI (with A{\footnotesize NSYS}) and FDS.}
\label{simple_panel_hot}
\end{figure}

\section{H profile column exposed to a localizes fire (\texorpdfstring{\textct{h\_profile}}{h\_profile})}

In this case a simple supported H-profile column is exposed to an adjacent hot panel ($T_s=1000$~\si{\degree C}). The steel column is 3~m tall (\textit{z} axis). The cross section is a 0.3~m (flange, parallel to \textit{x}) x 0.4~m (web, parallel to \textit{y}) with a 12.5~mm thick web and 16~mm thick flanges. The hot panel is parallel to the web and localized 30~cm from the web in \textit{x} direction (or 15~cm from the end of the flange).
The column is discretized with shell elements (FEM) to verify the use of the code to model shell structures. Each side of these plane elements will have a different thermal exposure, incident radiation and gas temperature around the surface. FDS2FTMI is designed to prescribe the correspondent heat flux at the shell elements top and bottom layers. Points A, B and C are located at the column middle span. A and C are located at the end of the flange (closest points to hot panel), B is located at the middle of the web.

The target of this qualitative verification exercise is to check the thermal gradient due to heat fluxes applied in top and bottom layers, the spatial temperature distribution and the thermomechanical behavior due to heating. The evolution of the surface temperature is presented in Fig.~\ref{h_profile1}. As the hot panel is parallel to the web, the heat flux incident to its surface will be higher than the heat flux incident to the flanges, so the surface temperature at A is higher than at B. Even with the high conductivity of steel is possible to observe a little gradient between the exposed and back surfaces at A and B. The spatial distribution also presented in Fig.~\ref{h_profile1} shows the expected profile. This profile is not exactly symmetric through column mid-span because of the energy reflected by the ground and the variation on the gas velocities around the column (which changes the heat transfer coefficient - $h$), so the temperature is lower at the top of the column.

\begin{figure}[ht]
\noindent
\parbox[c]{22em}{\includegraphics[width=3.0in]{scripts/SCRIPT_FIGURES/h_profile_ts}} 
\parbox[c]{11em}{\includegraphics[angle=270, trim=1cm 18cm 2cm 1.4cm, clip=true, origin=l,scale=0.30]{scripts/SCRIPT_FIGURES/h_profile_temp}} 
\parbox[c]{1em}{\includegraphics[angle=270, trim=1cm 18cm 2cm 1.4cm, clip=true, origin=l,scale=0.30]{scripts/SCRIPT_FIGURES/h_profile_temp3d}} 
\caption[The \textct{h\_profile} results]{Evolution of surface temperature and spatial distribution at 60 min.}
\label{h_profile1}
\end{figure}

This simple supported column is subjected to a vertical load of 325 kN, which correspond to 1/50 of the Euler's buckling critical load. Under these loading conditions, the column will be close to a uniform stress state and remain straight at the beginning of the simulation, then the horizontal displacements generated will be due to the temperature distribution. The evolution of horizontal displacements is presented in Fig.~\ref{h_profile2}. As expected, the horizontal displacements in \textit{x} axis are equal at A and C and represents the maximum achieved at the column, as can be observed at the lateral view. In \textit{y} direction the value of the displacement at these two points are the same, but as column is expanding and they are moving away from each other the signal of these displacements are opposite.

\begin{figure}[ht]
\noindent
\parbox[c]{19em}{\includegraphics[width=3.0in]{scripts/SCRIPT_FIGURES/h_profile_dx}} 
\parbox[c]{17.5em}{\includegraphics[width=2.785in, trim=1.5cm 0cm 0cm 0cm, clip=true]{scripts/SCRIPT_FIGURES/h_profile_dy}}
\parbox[c]{1em}{\includegraphics[angle=270, trim=1cm 18cm 3cm 1.4cm, clip=true, origin=l,scale=0.28]{scripts/SCRIPT_FIGURES/h_profile_disp_15x}} 
\caption[The \textct{h\_profile} results]{Evolution of horizontal displacements in \textit{x} and \textit{y} axis and a lateral view of the plane \textit{xz} (displacements amplified by x15).}
\label{h_profile2}
\end{figure}

\section{Complex Geometry Mapping (\texorpdfstring{\textct{h\_profile\_geom}}{h\_profile\_geom})}

This verification case extends the \textct{h\_profile} example to evaluate the new complex geometry mapping capabilities of FDS2FTMI. While the previous case used FDS orthogonal obstructions (\textct{OBST}) and \textct{BNDF} files, this case utilizes unstructured triangulated surfaces (\textct{GEOM}) and their associated \textct{.be} and \textct{.gcf} boundary files. The boundary mapping is achieved by calculating the closest-face centroid proximity between the FDS geometry faces and the FEM nodes.

To ensure the new \textct{GEOM} extraction methodology is accurate, a direct comparison is made against the \textct{OBST} mapping. Diagnostic parity plots evaluate the nodal extraction of the Adiabatic Surface Temperature ($T_{\hbox{\tiny AST}}$) and the convective heat transfer coefficient ($h$). Figure~\ref{geom_parity_plots} displays these parity plots, demonstrating a 1:1 perfect match between the mapped boundary conditions.

\begin{figure}[ht]
\noindent
\begin{tabular*}{\textwidth}{l@{\extracolsep{\fill}}r}
\includegraphics[width=3.0in]{scripts/SCRIPT_FIGURES/t_ast_mapping_parity.pdf} &
\includegraphics[width=3.0in]{scripts/SCRIPT_FIGURES/h_mapping_parity.pdf}
\end{tabular*}
\caption[The \textct{h\_profile\_geom} parity results]{Diagnostic parity plots comparing mapped boundary conditions ($T_{\hbox{\tiny AST}}$ and $h$) between \textct{OBST} and \textct{GEOM} extraction methods.}
\label{geom_parity_plots}
\end{figure}

Furthermore, the transient thermal evolution and the resulting thermomechanical displacements are tracked and compared. Figure~\ref{geom_thermal} shows the thermal history comparison for points A and B, confirming that the \textct{OBST} overlaps with the \textct{GEOM} results.

\begin{figure}[ht]
\centering
\includegraphics[width=4.0in]{scripts/SCRIPT_FIGURES/h_profile_geom_thermal.pdf}
\caption[The \textct{h\_profile\_geom} thermal results]{Comparison of the temperature evolution between the traditional orthogonal mesh and the unstructured complex geometry.}
\label{geom_thermal}
\end{figure}

The structural response due to these thermal profiles is verified by tracking the horizontal displacements. Figure~\ref{geom_displacements} displays the displacement tracking in both the $x$ and $y$ axes. The overlapping displacement curves confirm that the unstructured \textct{GEOM} mapping yields equivalent thermomechanical behavior.

\begin{figure}[ht]
\noindent
\begin{tabular*}{\textwidth}{l@{\extracolsep{\fill}}r}
\includegraphics[width=3.0in]{scripts/SCRIPT_FIGURES/h_profile_geom_disp_x.pdf} &
\includegraphics[width=3.0in]{scripts/SCRIPT_FIGURES/h_profile_geom_disp_y.pdf}
\end{tabular*}
\caption[The \textct{h\_profile\_geom} displacement results]{Comparison of horizontal displacements in the $x$ and $y$ axes between traditional and complex geometry models.}
\label{geom_displacements}
\end{figure}

Quantitative error metrics for maximum temperature, maximum $x$-displacement, and maximum $y$-displacement are calculated and evaluated against strict tolerances (0.025 for temperature, and 0.01 for displacements). These metrics are automatically appended to the overall verification statistics.

\section{Summary of Verification Results}

This section summarizes the results from Verification cases with the objective of evaluate the output predicted by FDS2FTMI code. The presented results can help identify errors in the source code.

\begin{longtable}[c]{l c c c c c c }
\hline
Case Name & Expected & Predicted & Type of Error & Error & Error     & Within    \\
          & Metric   & Metric    &               &       & Tolerance & Tolerance \\ \hline \hline
\endfirsthead
\hline
Case Name & Expected & Predicted & Type of Error & Error & Error     & Within    \\
          & Metric   & Metric    &               &       & Tolerance & Tolerance \\ \hline \hline
\endhead
\hline
\endfoot
\hline
\endlastfoot
\input{scripts/verification_statistics.tex}
\end{longtable}

\printbibliography
\end{document}
